% Source: Weinberg GR, physical PDF pp. 95--98
%         (printed pp. 70--73).
% Coverage: Section 3.2, Gravitational Forces; equations
%           (3.2.1)--(3.2.14).
% Figures/tables/footnotes: none.
% Status: source-reviewed and compile-clean.
% Uncertainties: none.

\subsection{Gravitational Forces}\label{sec:3.2}

Consider a particle moving freely under the influence of purely
gravitational forces. According to the Principle of Equivalence, there is a
freely falling coordinate system $\xi^\rho$ in which its equation of motion
is that of a straight line in spacetime, that is,
\begin{equation}
  \frac{\dd^2\xi^\rho}{\dd\tau^2}=0,
  \tag{3.2.1}\label{eq:3.2.1}
\end{equation}
with $\dd\tau$ the proper time,
\begin{equation}
  \dd\tau^2=-\eta_{\rho\sigma}\,\dd\xi^\rho\dd\xi^\sigma.
  \tag{3.2.2}\label{eq:3.2.2}
\end{equation}
[Compare Eqs.~(2.3.1) and (2.1.4).] Now suppose that we use any other
coordinate system $x^\mu$, which may be a Cartesian coordinate system at
rest in the laboratory, but also may be curvilinear, accelerated, rotating,
or what we will. The freely falling coordinates $\xi^\rho$ are functions of
the $x^\mu$, and Eq.~\eqref{eq:3.2.1} becomes
\[
  0=\frac{\dd}{\dd\tau}
  \left(\frac{\partial\xi^\rho}{\partial x^\mu}
  \frac{\dd x^\mu}{\dd\tau}\right)
  =\frac{\partial\xi^\rho}{\partial x^\mu}
  \frac{\dd^2x^\mu}{\dd\tau^2}
  +\frac{\partial^2\xi^\rho}{\partial x^\mu\partial x^\nu}
  \frac{\dd x^\mu}{\dd\tau}\frac{\dd x^\nu}{\dd\tau}.
\]
Multiply this by $\partial x^\lambda/\partial\xi^\rho$, and use the
familiar product rule
\[
  \frac{\partial\xi^\rho}{\partial x^\mu}
  \frac{\partial x^\lambda}{\partial\xi^\rho}
  =\delta^\lambda_\mu.
\]
This gives the equation of motion
\begin{equation}
  \frac{\dd^2x^\lambda}{\dd\tau^2}
  +\Gamma^\lambda{}_{\mu\nu}
  \frac{\dd x^\mu}{\dd\tau}\frac{\dd x^\nu}{\dd\tau}=0,
  \tag{3.2.3}\label{eq:3.2.3}
\end{equation}
where $\Gamma^\lambda{}_{\mu\nu}$ is the \emph{affine connection}, defined by
\begin{equation}
  \Gamma^\lambda{}_{\mu\nu}
  \equiv
  \frac{\partial x^\lambda}{\partial\xi^\rho}
  \frac{\partial^2\xi^\rho}{\partial x^\mu\partial x^\nu}.
  \tag{3.2.4}\label{eq:3.2.4}
\end{equation}

The proper time \eqref{eq:3.2.2} may also be expressed in an arbitrary
coordinate system,
\begin{equation}
  \dd\tau^2
  =-\eta_{\rho\sigma}
  \frac{\partial\xi^\rho}{\partial x^\mu}\dd x^\mu
  \frac{\partial\xi^\sigma}{\partial x^\nu}\dd x^\nu,
  \tag{3.2.5}\label{eq:3.2.5}
\end{equation}
or
\begin{equation}
  \dd\tau^2=-g_{\mu\nu}\,\dd x^\mu\dd x^\nu,
  \tag{3.2.6}\label{eq:3.2.6}
\end{equation}
where $g_{\mu\nu}$ is the \emph{metric tensor}, defined by
\begin{equation}
  g_{\mu\nu}\equiv
  \frac{\partial\xi^\rho}{\partial x^\mu}
  \frac{\partial\xi^\sigma}{\partial x^\nu}\eta_{\rho\sigma}.
  \tag{3.2.7}\label{eq:3.2.7}
\end{equation}

For a photon or a neutrino the equation of motion in a freely falling system
is the same as \eqref{eq:3.2.1}, except that the independent variable cannot
be taken as the proper time \eqref{eq:3.2.2}, because for massless particles
the right-hand side of \eqref{eq:3.2.2} vanishes. Instead of $\tau$ we can
use $\sigma\equiv\xi^0$, so that \eqref{eq:3.2.1} and \eqref{eq:3.2.2}
become
\[
  \frac{\dd^2\xi^\rho}{\dd\sigma^2}=0,
  \qquad
  0=-\eta_{\rho\sigma}
  \frac{\dd\xi^\rho}{\dd\sigma}\frac{\dd\xi^\sigma}{\dd\sigma}.
\]
Following the same reasoning as before, we find that the equation of motion
in an arbitrary gravitational field and an arbitrary coordinate system is
\begin{equation}
  \frac{\dd^2x^\mu}{\dd\sigma^2}
  +\Gamma^\mu{}_{\nu\lambda}
  \frac{\dd x^\nu}{\dd\sigma}\frac{\dd x^\lambda}{\dd\sigma}=0,
  \tag{3.2.8}\label{eq:3.2.8}
\end{equation}
\begin{equation}
  0=-g_{\mu\nu}
  \frac{\dd x^\mu}{\dd\sigma}\frac{\dd x^\nu}{\dd\sigma},
  \tag{3.2.9}\label{eq:3.2.9}
\end{equation}
with $\Gamma^\mu{}_{\nu\lambda}$ and $g_{\mu\nu}$ given as before by
\eqref{eq:3.2.4} and \eqref{eq:3.2.7}.

Incidentally, in both \eqref{eq:3.2.3} and \eqref{eq:3.2.8} we do not need
to know what $\tau$ and $\sigma$ are in order to find the motion of our
particle, for these equations when solved give $x^\mu(\tau)$ or
$x^\mu(\sigma)$, and $\tau$ or $\sigma$ can be eliminated to give
$\mathbf x(t)$. The purpose of \eqref{eq:3.2.6} is to tell us how to compute
the proper time, whereas the purpose of \eqref{eq:3.2.9} is to impose initial
conditions appropriate to a massless particle. In particular,
Eq.~\eqref{eq:3.2.9} tells us that the time $\dd t$ for a photon to travel a
distance $\dd\mathbf x$ is determined by the quadratic equation
\[
  0=g_{00}\,\dd t^2+2g_{i0}\,\dd x^i\dd t
  +g_{ij}\,\dd x^i\dd x^j,
\]
with $i$ and $j$ summed over the values $1,2,3$. The solution is
\begin{equation}
  \dd t=\frac{1}{g_{00}}\left[
  -g_{i0}\,\dd x^i
  -\left\{(g_{i0}g_{j0}-g_{ij}g_{00})
  \dd x^i\dd x^j\right\}^{1/2}\right],
  \tag{3.2.10}\label{eq:3.2.10}
\end{equation}
and the time required for light to travel along any path may be calculated
by integrating $\dd t$ along the path.

The values of the metric tensor $g_{\mu\nu}$ and the affine connection
$\Gamma^\lambda{}_{\mu\nu}$ at a point $X$ in an arbitrary coordinate
system $x^\mu$ provide enough information to determine the locally inertial
coordinates $\xi^\rho(x)$ in a neighborhood of $X$. First, we multiply
Eq.~\eqref{eq:3.2.4} by $\partial\xi^\sigma/\partial x^\lambda$ and use the
product rule
\[
  \frac{\partial\xi^\sigma}{\partial x^\lambda}
  \frac{\partial x^\lambda}{\partial\xi^\rho}
  =\delta^\sigma_\rho,
\]
thereby obtaining the differential equations for $\xi^\rho$,
\begin{equation}
  \frac{\partial^2\xi^\rho}{\partial x^\mu\partial x^\nu}
  =\Gamma^\lambda{}_{\mu\nu}
  \frac{\partial\xi^\rho}{\partial x^\lambda}.
  \tag{3.2.11}\label{eq:3.2.11}
\end{equation}
The solution is
\begin{align}
  \xi^\rho(x)={}&a^\rho+b^\rho{}_\mu(x^\mu-X^\mu)
  \notag\\
  &+\frac12 b^\rho{}_\lambda\Gamma^\lambda{}_{\mu\nu}(X)
  (x^\mu-X^\mu)(x^\nu-X^\nu)+\cdots,
  \tag{3.2.12}\label{eq:3.2.12}
\end{align}
where
\begin{equation}
  a^\rho=\xi^\rho(X),
  \qquad
  b^\rho{}_\lambda=
  \frac{\partial\xi^\rho(X)}{\partial X^\lambda}.
  \tag{3.2.13}\label{eq:3.2.13}
\end{equation}
From Eq.~\eqref{eq:3.2.7} we also learn that
\begin{equation}
  \eta_{\rho\sigma}b^\rho{}_\mu b^\sigma{}_\nu=g_{\mu\nu}(X).
  \tag{3.2.14}\label{eq:3.2.14}
\end{equation}

Thus, given $\Gamma^\lambda{}_{\mu\nu}$ and $g_{\mu\nu}$ at $X$, the locally
inertial coordinates $\xi^\rho$ are determined to order $(x-X)^2$, except
for the ambiguity in the constants $a^\rho$ and $b^\rho{}_\lambda$. The
$b^\rho{}_\lambda$ are determined by Eq.~\eqref{eq:3.2.13} up to a Lorentz
transformation
$b^\rho{}_\mu\to\Lambda^\rho{}_\sigma b^\sigma{}_\mu$, so the ambiguity in
the solution for $\xi^\rho(x)$ just reflects the fact that if $\xi^\rho$ are
locally inertial coordinates, then so are
$\Lambda^\rho{}_\sigma\xi^\sigma+c^\rho$. Hence, since
$\Gamma^\lambda{}_{\mu\nu}$ and $g_{\mu\nu}$ determine the locally inertial
coordinates up to an inhomogeneous Lorentz transformation, and since the
gravitational field can have no effects in a locally inertial coordinate
system, we should not be surprised to find that all effects of gravitation
are comprised in $\Gamma^\lambda{}_{\mu\nu}$ and $g_{\mu\nu}$. Note,
however, that \eqref{eq:3.2.12} satisfies \eqref{eq:3.2.11} only at the point
$x=X$; in order for it to be possible to solve \eqref{eq:3.2.11} for all
$x$, it is necessary for the derivatives of the affine connection to
satisfy certain symmetry conditions, to be discussed in Chapter~5.
